\chapter{Reproducibility, Defects and Research Integrity}\label{ch:reproducibility}

\needswork{STATUS: DRAFTED. This chapter is unusual for a thesis and deliberate. It reports the
errors found in this project's own results. Examiners should read it as the methodological
contribution it is intended to be, not as an admission appended to a limitations section.}

\section{Why this chapter exists}

Most theses report what was found. This one also reports \emph{what was found to be wrong, and how}.
The justification is not confession but transferability: the defects catalogued below are not
peculiar to this work, and several of them survived safeguards that are widely considered
sufficient --- version control, a green test suite, seeded runs, committed raw data.

Ten results in this project moved after they were first recorded. Only four were catchable by
running more seeds. Two were \emph{protected by passing tests}. Three were contradictions between a
paper and the artifact that generated it.

\section{The defect taxonomy}

\begin{center}
\begin{tabular}{p{2.2cm}p{5.6cm}p{5.2cm}}
\toprule
class & mechanism & instance in this work \\
\midrule
\textbf{C1} & small-sample mean lands the wrong side of a threshold & four headline numbers; a
capacity value $5 \rightarrow 3$ \\
\textbf{C2} & an unverified constant on the measurement path & the $2.4$\,GHz basic-rate
saturation; the frame header treated as a constant \\
\textbf{C3} & a threshold applied to a mean rather than a distribution & the capacity criterion;
the utilisation crossing \\
\textbf{C4} & a configuration change that perturbs the random realisation & simulator stream
displacement under mobility \\
\textbf{C5} & a claim wider than the experiment supports & two-node hardware described as
validating contention \\
\bottomrule
\end{tabular}
\end{center}

\begin{remark}[The load-bearing observation]
\textbf{More seeds fixes only C1.} C2--C5 are immune to sample size. ``We ran thirty seeds'' is
therefore not a statement of correctness, and the discipline installed after the C1 failures gave
no protection at all against the C2 failures found later.
\end{remark}

\section{An automated staleness gate, and what it does not cover}

Every model-derived artifact is re-derived from current code on every run and compared row by row;
a single differing value fails the build. This makes silent staleness of \emph{derived} data
impossible.

It does not make staleness impossible. Two limits were found the hard way:

\begin{itemize}
\item \textbf{Scope.} Every drift guard originally parsed the paper. The specification documents
were outside that scope, so corrections propagated into the paper and not into the documents for
three weeks --- including a table built verbatim on data the project had already purged and
superseded.
\item \textbf{Direction.} Retracting a finding in a register does not remove it from the prose. One
retracted claim was found alive in the paper months later, one hundred lines below a sentence
stating the opposite, and later again in the open-items register. \textbf{When you retract, grep the
whole repository.}
\end{itemize}

\section{Practices that earned their place}

\begin{description}
\item[Pre-registration, committed data-free.] Predictions are committed as their own commit before
data exist, so the ordering is third-party checkable. This cost the project a claim once: an
assertion of pre-registration was \emph{withdrawn} rather than reconstructed, because writing an
expectations file after results are known manufactures evidence.

\item[Stating the expected value before measuring.] This is what caught the $2.4$\,GHz
basic-rate artifact, which repeatability would never have caught.

\item[Retaining a discarded model as a labelled failure curve.] The naive broadcast reduction
remains in the codebase, clearly marked, so the $16\times$ error it produces can be shown rather
than described.

\item[Keeping per-seed raw data, not summaries.] This allowed validation bands to be re-derived at
thirty seeds without re-simulating.

\item[An independent second implementation.] The slot-exact simulator, written before the published
model was found, shares none of its assumptions and is what turns agreement into evidence.
\end{description}

\section{Two defects worth stating individually}

\subsection{A test that asserted a bug}

A utilisation function returned exactly zero for a single node, with a unit test asserting that a
lone sender never contends. The function conflated \emph{contention} with \emph{utilisation}: a lone
sender does not collide but still occupies airtime, and returning zero declared a single node able
to carry unbounded traffic. The defect was harmless until another result began reporting
single-node cases --- \textbf{and it was protected by a passing test that had encoded the wrong
behaviour as intended behaviour}.

\subsection{A criterion that could not fail}

The project pre-registered a success criterion requiring at least a $40\%$ reduction in
authentication bytes, and committed it to version control two days before the result existed. The
ordering is genuine and verifiable.

The quantity it tests, however, reduces exactly to $1 - 1/b$. The threshold is therefore
$b \geq 2$ --- independent of encoding, scheme, placement, header size and signature size. The
companion condition on verifiability is satisfied by construction with zero margin. \textbf{No
configuration that batches at all could have failed the test.}

The pre-registration is real and worth keeping; what had to go was the implication that it
constituted a risk the design might not have survived.

\section{What a reader should take from this}

Three times in this project a claim survived because two facts were each individually recorded and
never placed side by side. A register that stores facts separately does not compose them.
\textbf{Only re-derivation does.}

\needswork{TO ADD: a comparison with the reproducibility literature --- the simulation-credibility
lineage in particular --- so this chapter reads as a contribution to that conversation rather than
as a project post-mortem. The companion methods paper does this; the material needs porting.}
